\documentclass[10pt]{beamer}
\usetheme{metropolis}

\usepackage{amsmath,amssymb}
\usepackage{mathtools}
\usepackage{xeCJK}
\setCJKmainfont{Adobe 楷体 Std}
\usepackage{tikz-cd}

% 1. 开启标题背景色填充
\metroset{block=fill}

% theorem environments
\newtheorem{proposition}{Proposition}
\newtheorem{exercise}{Exercise}

% notation
\newcommand{\R}{\mathbb{R}}
\newcommand{\N}{\mathbb{N}}
\newcommand{\F}{\mathcal{F}}
\newcommand{\supp}{\operatorname{supp}}

\title{Ch1\quad Manifolds (continued)}
\subtitle{Second Countability \& Partition of Unity}
\date{September 22, 2026}
\author{LXJ}

\begin{document}

\maketitle

%====================================================================
\section{Paracompactness}

\begin{frame}{Covers, refinements and local finiteness}
	\begin{definition}[Cover, subcover, refinement]
		A collection $\{U_\alpha\}$ is a \emph{cover} of $W \subset M$ if
		$W \subset \bigcup_\alpha U_\alpha$. It is an \emph{open cover} if each
		$U_\alpha$ is open. A subcollection of $\{U_\alpha\}$ which still
		covers is called a \emph{subcover}. A \emph{refinement} $\{V_\beta\}$
		of the cover $\{U_\alpha\}$ is a cover such that
		\[
			\forall\, \beta,\ \exists\, \alpha \quad \text{s.t.}\quad
			V_\beta \subset U_\alpha .
		\]
	\end{definition}
	\begin{definition}[Locally finite]
		A collection $\{A_\alpha\}$ of subsets of $M$ is \emph{locally finite}
		if whenever $m \in M$, $\exists$ a neighborhood $W_m$ of $m$ s.t.\
		\[
			W_m \cap A_\alpha \neq \varnothing
			\quad\text{for only finitely many } \alpha .
		\]
	\end{definition}
\end{frame}

\begin{frame}{Paracompactness}
	\begin{definition}[Paracompact]
		A topological space is \emph{paracompact} if every open cover has a
		\emph{locally finite refinement}.
	\end{definition}
	\begin{proposition}
		Let $X$ be a topological space which is
		\begin{itemize}
			\item \emph{locally compact} (each point has a compact
			      neighborhood),
			\item \emph{Hausdorff}, and
			\item \emph{second countable}.
		\end{itemize}
		Then $X$ is paracompact.
	\end{proposition}
	\begin{alertblock}{Corollary}
		Every differentiable manifold is paracompact.
	\end{alertblock}
\end{frame}

\begin{frame}{Proof sketch of the proposition}
	\begin{proof}[Sketch]
		\begin{enumerate}
			\item \textbf{Exhaustion.} $X$ has a countable basis
			      $\{B_n\}_{n \in \N^+}$ of open sets with compact closure.
			      Inductively build $G_1 = B_1$,
			      $G_{k+1} = \bigcup_{n=1}^{j_{k+1}} B_n$ with $j_{k+1}$
			      smallest s.t.\ $\overline{G_k} \subset
				      \bigcup_{n=1}^{j_{k+1}} B_n$. Then
			      \[
				      G_1 \subset G_2 \subset \cdots \subset G_i \subset G_{i+1}
				      \subset \cdots, \qquad \textstyle\bigcup_k G_k = X .
			      \]
			\item \textbf{Refinement.} Given an open cover $\{U_\alpha\}$:
			      $\overline{G_i} \setminus G_{i-1}$ is compact, covered by
			      finitely many $U_\alpha \cap (G_{i+1} \setminus
				      \overline{G_{i-2}})$; similarly cover $\overline{G_2}$ by
			      finitely many $U_\alpha \cap G_3$. Collecting these gives
			      a countable refinement $\{V_n\}$.
			\item \textbf{Local finiteness.} Each $p \in X$ lies in at most two
			      ``rings'' $\overline{G_i} \setminus G_{i-1}$, each meeting
			      only finitely many $V_n$.
		\end{enumerate}
	\end{proof}
\end{frame}

\begin{frame}{An exercise hidden in the proof}
	\begin{exercise}
		If $X$ is second countable, locally compact and Hausdorff, then $X$
		has a countable basis consisting of open sets with compact closure.
	\end{exercise}
	\begin{block}{Remark}
		The picture behind the construction:
		\[
			\overline{G_i} \setminus G_{i-1}
			\ \subset\
			G_{i+1} \setminus \overline{G_{i-2}}
			\qquad (\text{compact ring inside open ring}).
		\]
		Each $V_n$ has compact closure, being contained in
		$\overline{G_{i+1}} \setminus G_{i-2}$ or $\overline{G_3}$.
	\end{block}
\end{frame}

%====================================================================
\section{Partition of Unity}

\begin{frame}{Partition of unity}
	\begin{definition}[Partition of unity]
		A \emph{partition of unity} on $M$ is a collection
		$\{\varphi_i \mid i \in I\}$ of $C^\infty$ functions on $M$ s.t.
		\begin{enumerate}
			\item[(a)] the collection $\{\supp \varphi_i\}_{i \in I}$ is
			      \emph{locally finite};
			\item[(b)] $\displaystyle\sum_{i \in I} \varphi_i = 1$ on $M$, and
			      $\varphi_i(p) \geq 0$ for all $p \in M$ and $i \in I$.
		\end{enumerate}
	\end{definition}
	\begin{definition}[Subordination]
		A partition of unity $\{\varphi_i \mid i \in I\}$ is
		\emph{subordinated to the cover} $\{U_\alpha\}_{\alpha \in A}$ if
		\[
			\forall\, i \in I,\ \exists\, \alpha \in A \quad \text{s.t.}\quad
			\supp \varphi_i \subset U_\alpha .
		\]
		It is subordinated to $\{U_i\}_{i \in I}$ \emph{with the same index}
		if $\supp \varphi_i \subset U_i$ for each $i \in I$.
	\end{definition}
\end{frame}

\begin{frame}{Bump functions on $\R^d$}
	\begin{proposition}[Existence of bump functions]
		There exists a non-negative $C^\infty$ function $\varphi$ on $\R^d$
		which equals $1$ on the closed cube $\overline{C(1)}$ and zero on the
		complement of the open cube $C(2)$.
		\par\smallskip
		Similarly, $\exists$ a non-negative $C^\infty$ function $\varphi$ on
		$\R^d$ which equals $1$ on the closed ball $\overline{B_0(1)}$ and
		zero outside of the open ball $B_0(2)$.
	\end{proposition}
	\begin{proof}[Sketch]
		\begin{enumerate}
			\item $f(t) = \begin{cases} e^{-1/t^2}, & t > 0 \\ 0, & t \leq 0 \end{cases}$
			      is $C^\infty$, and $f(t) \to 1^{-}$ as $t \to +\infty$.
			\item $g(t) = \dfrac{f(t)}{f(t) + f(1-t)}$ is $C^\infty$
			      (denominator never zero), with $g(t) \equiv 0$ ($t \leq 0$)
			      and $g(t) \equiv 1$ ($t \geq 1$).
			\item $h(t) = g(t+2)\, g(2-t)$: then $h \equiv 1$ on $[-1, 1]$ and
			      $h \equiv 0$ outside $[-2, 2]$.
			\item General case: $\varphi = (h \circ r_1) \cdots (h \circ r_d)$
			      for the cube, or $\varphi = h\bigl(\|x\|^2\bigr)$ for the ball.
		\end{enumerate}
	\end{proof}
\end{frame}

\end{document}
